\section{Итоговые проекты: самостоятельные миссии}
\label{sec:python-final}

\begin{infobox}[Цели главы]
\begin{itemize}
    \item Соединить маршрут, камеру, распознавание, наведение и посадку в одной программе.
    \item Для каждого этапа заранее определить наблюдаемый результат и условие перехода к следующему.
    \item Проверять не только успех, но и ограниченные неуспешные исходы: цели нет, цель потеряна, карта другая.
\end{itemize}
\end{infobox}

\begin{missionbox}[Миссии симулятора к главе]
Язык \ui{Python}, панель \ui{Обучение}, \textbf{100--108 «Самостоятельные миссии»}: карта из нескольких меток, поиск по секторам, лесная экспедиция, осмотр крыш, контрольные точки арены, инспекция промышленной зоны, неожиданная потеря цели, незнакомая карта, итоговая экспедиция. Новых команд в них нет: всё нужное --- в главах~\ref{sec:python-sdk} и~\ref{sec:python-vision}.
\end{missionbox}

\begin{rulebox}[Как оформлять проект]
Программа собирается из функций прошлых глав: \verb|wait_state|, \verb|go| (точка с ожиданием и тайм-аутом), \verb|read_target| (кадр и углы метки), \verb|observation| (ошибка метки), \verb|align| (наведение). У каждого этапа --- своё сообщение в журнале и своё условие перехода; любой поиск и любое ожидание ограничены; при неудаче программа сообщает причину, обнуляет скорость и сажает дрон. Номер метки в коде не превращают в координаты: где метка, говорит только кадр.
\end{rulebox}

\begin{infobox}[Общие критерии оценки]
\begin{enumerate}
    \item \textbf{Результат на сцене:} дрон делает то, что требует проект, и садится; при неудаче --- безопасно садится с понятным сообщением.
    \item \textbf{Журнал:} по выводу программы восстанавливается порядок этапов, прочитанные ID и причина остановки.
    \item \textbf{Ограниченность:} нет бесконечных ожиданий, движения по старому кадру и посадки «по цели» без подтверждения.
    \item \textbf{Читаемость:} код разбит на функции, параметры (ID, высоты, допуски) вынесены в переменные.
\end{enumerate}
\end{infobox}

\subsection{Метки как карта}

Несколько меток превращают полигон в карту с контрольными местами. Назначение метки хранят в словаре Python: \verb|names = {10: "северо-запад", 11: "юго-восток", 12: "цель"}|, где ключ --- ID, значение --- смысл. Проверка \verb|marker_id in names| отвечает, есть ли такой ID в таблице.

\begin{taskbox}[Карта из нескольких меток]
С высоты 3~м над стартом прочитайте все видимые метки и для каждой напечатайте «ID: назначение» по таблице \verb|names|. Метки, которых в таблице нет, печатайте как «ID: неизвестна». После отчёта сядьте.
\expected{в журнале три строки: \texttt{10:северо-запад}, \texttt{11:юго-восток}, \texttt{12:цель} (в любом порядке), напечатаны только реально прочитанные ID; дрон садится на старте.}
\mission{Python}{100}
\end{taskbox}

\subsection{Поиск по секторам}

Большую область делят на участки, которые камера видит целиком: с высоты 3~м нижняя камера видит около $\pm2$~м по Y, поэтому сектора ставят через 4~м. Список секторов конечен, каждый посещается один раз, после находки поиск прекращается.

\begin{taskbox}[Поиск по секторам]
Найдите метку 12, проверяя сектора $Y = 0, 4, 8$ при $X = 0$ на высоте 3~м. В каждом секторе печатайте его номер; после находки --- «Найдена: 12» и прекратите поиск. Если метки нет ни в одном секторе, напечатайте «Поиск истёк». Вернитесь на старт и сядьте.
\expected{в сцене миссии проверяются два пустых сектора и сектор с целью, после находки дрон сразу летит домой; с несуществующим ID посещаются все три сектора и печатается \texttt{Поиск истёк}.}
\mission{Python}{101}
\end{taskbox}

\subsection{Экспедиции по полигонам}

Маршрут доставляет к месту, камера подтверждает цель: это две разные проверки. Высоту и путь выбирают по местности (раздел~\ref{sec:python-terrain}), а кадр снимают уже после прибытия и паузы, когда корпус успокоился.

\begin{taskbox}[Лесная экспедиция и крыши]
В лесу пролетите центральный коридор ($X = 0$) до точки осмотра $(0, 8, 3)$, прочитайте ID~12, вернитесь и сядьте. В спальном районе наберите над стартом высоту 9~м, осмотрите крыши домов в $(9, -7)$ и $(-9, -7)$, напечатайте прочитанные ID и вернитесь на старт по той же высоте.
\expected{в лесу дрон не касается деревьев и печатает \texttt{Прочитан: 12}; над домами печатаются ID 102 и 101, высота везде выше крыш с запасом.}
\mission{Python}{102--103}
\end{taskbox}

\begin{taskbox}[Контрольные точки арены]
На Геоскан Арене подлетите на высоте 6~м к карте ArUco в $(-9, 10.5)$, снизьтесь до 4~м и осмотрите её из двух точек: над $(-9, 10.5)$ и со сдвигом на 0.5~м по $+X$. В каждой точке напечатайте координаты прибытия и прочитанные ID, затем вернитесь на старт через высоту 6~м и сядьте.
\expected{из обеих точек в журнале ID 10 и 11 рядом с координатами точки; во втором кадре метки смещены влево относительно первого; посадка на старте.}
\mission{Python}{104}
\end{taskbox}

\begin{taskbox}[Инспекция промышленной зоны]
Задайте участки списком пар \verb|[(8, 12), (16, 13)]| --- координата Y участка на линии $X = 0$ и ID, который должен там быть. На каждом участке ищите именно его ID; если цель участка не найдена, прервите инспекцию с сообщением и сядьте.
\expected{в журнале по порядку \texttt{Прочитан: 12} и \texttt{Прочитан: 13}; если поменять ID в списке местами, программа останавливается на первом участке с сообщением «Цель участка не найдена».}
\mission{Python}{105}
\end{taskbox}

\subsection{Потеря цели и незнакомая карта}

Надёжная программа знает, что делать, когда всё пошло не так: при потере цели --- обнулить скорость и вести ограниченное восстановление; на новой карте --- не полагаться на запомненные координаты.

\begin{taskbox}[Неожиданная потеря цели]
Найдите метку 12 над стартом, отлетите на $(0, 5, 3)$, где её не видно, и восстановите цель не больше чем за три попытки: на каждой обнулите скорость и снимите кадр, на второй вернитесь к старту. При обнаружении напечатайте «Восстановлена» и прекратите поиск, иначе --- «Поиск истёк». Сядьте.
\expected{журнал: «Потеряна», затем «Восстановлена» после возврата; если убрать возврат --- три строки «Потеряна» и \texttt{Поиск истёк}.}
\mission{Python}{106}
\end{taskbox}

\begin{taskbox}[Незнакомая карта]
Запустите программу посадки на метку (\texttt{12\_01\_land\_on\_marker.py}) без изменений на смешанной карте миссии 107, где метка 12 стоит в другом месте. Координаты метки в программу не записывайте. Напечатайте точку посадки и сравните её с $(0.6, 0.3)$.
\expected{дрон садится на метку 12 там, где она стоит на этой карте; напечатанная точка посадки заметно отличается от $(0.6, 0.3)$.}
\mission{Python}{107}
\end{taskbox}

\subsection{Итоговая экспедиция}

\begin{taskbox}[Итоговая экспедиция]
Соберите полную миссию с параметром \verb|target_id|: подготовка и взлёт с подтверждением, поиск метки по секторам $Y = 0, 1, -1$, намеренная потеря (отлёт на $(0, 5, 3)$), восстановление через точки $Y = 4$ и $0$, наведение \verb|align|, посадка на метку и отчёт. Запустите дважды: с \verb|target_id = 12| и с \verb|target_id = 249|, которого на карте нет.

Критерии проверки в симуляторе:
\begin{enumerate}
    \item С ID~12: в журнале по порядку «Цель миссии: 12», находка, «Потеряна», «Восстановлена», центр подтверждён, дрон сел на метку.
    \item С ID~249: поиск обходит три сектора, печатает «Поиск истёк», дрон возвращается на старт и садится без ложного сообщения об успехе.
    \item Во время потери цели нет ненулевых команд скорости; все ожидания ограничены по времени или числу попыток.
    \item Каждый переход между этапами объяснён в журнале свежим наблюдением, которое его разрешило.
\end{enumerate}
\expected{оба запуска завершаются посадкой, отчёт первого --- успех с восстановлением, второго --- тайм-аут поиска.}
\mission{Python}{108}
\end{taskbox}

\subsection{Проекты для настоящего дрона}

\begin{taskbox}[Автономный инспектор]
Только на дроне. Облетите четыре «стеллажа» --- маркеры ID 0, 1, 2, 3 по углам комнаты. У каждого найдите маркер поворотом на месте (ограниченный поиск), наведитесь на него и зафиксируйте в журнале; затем вернитесь на старт и сядьте.
\expected{в журнале четыре подтверждённых ID в порядке обхода; если маркер не найден за отведённое время, программа пропускает стеллаж с сообщением и продолжает маршрут.}
\end{taskbox}

\begin{taskbox}[Следуй за мной]
Только на дроне. Удерживайте человека в центре кадра на безопасном расстоянии: поворот \verb|yaw_rate| пропорционален смещению центра человека (MediaPipe, раздел~\ref{sec:human-tracking}); если плечи занимают больше половины ширины кадра --- отлёт назад, если меньше четверти --- вперёд. Скорости ограничены, при потере человека дольше 3~с --- зависание и посадка.
\expected{дрон поворачивается вслед за человеком и держит дистанцию, не приближаясь ближе 1.5~м; когда человек уходит из кадра, дрон останавливается и садится.}
\end{taskbox}

\begin{taskbox}[Гонка по воротам]
Только на дроне (или на гоночном полигоне симулятора без камеры). Пролетите через ворота, отмеченные ArUco-маркерами, как можно быстрее: для каждых ворот --- подводящая точка, центр проёма, выходная точка. Усложнение: ворота стоят на разной высоте.
\expected{дрон проходит каждые ворота через середину проёма, не задевая рамку; время прохождения печатается в конце.}
\mission{Python}{61}
\end{taskbox}
